(a) O grupo adimensional proposto é
\[
\Pi = \frac{P}{E L^2},
\]
cuja adimensionalidade se verifica substituindo as dimensões ($[P]=MLT^{-2}$, $[E]=ML^{-1}T^{-2}$, $[L^2]=L^2$):
\[
\left[\frac{P}{EL^2}\right] = \frac{MLT^{-2}}{ML^{-1}T^{-2}\cdot L^2} = \frac{MLT^{-2}}{MLT^{-2}} = 1.\quad\checkmark
\]

(b) Pela igualdade de grupos entre modelo (m) e protótipo (p):
\[
\frac{P_m}{E_m L_m^2}=\frac{P_p}{E_p L_p^2}
\implies
P_m=P_p\,\frac{E_m}{E_p}\left(\frac{L_m}{L_p}\right)^2.
\]
Dados: $P_p=9000$ N, $L_m=0{,}20$ m, $L_p=3{,}6$ m.

Módulos de elasticidade típicos: a\c{c}o $E_m \approx 200\,\mathrm{GPa}$, madeira $E_p \approx 12\,\mathrm{GPa}$.

Substituindo:
\[
P_m=9000\times\frac{200}{12}\times\left(\frac{0{,}20}{3{,}6}\right)^2
= 9000\times 16{,}7\times 0{,}00309
\approx 463\,\mathrm{N}.
\]

Portanto, deve-se aplicar uma carga de aproximadamente $\mathbf{463\,N}$ ao modelo de a\c{c}o para reproduzir as condi\c{c}\~oes do protótipo de madeira.

